\chapter{三角 Fourier 级数}

\begin{solution}{fourier-circle-real-coefficients}
若
\[
 a_k=\frac1\pi\int_{-\pi}^{\pi}f(x)\cos kx\,dx,\qquad
 b_k=\frac1\pi\int_{-\pi}^{\pi}f(x)\sin kx\,dx,
\]
则 \(k\ge1\) 时
\(\widehat f(k)=(a_k-ib_k)/2\)、\(\widehat f(-k)=(a_k+ib_k)/2\)，
且 \(\widehat f(0)=a_0/2\)。若 \(f\) 实值，对积分取共轭便得
\(\overline{\widehat f(k)}=\widehat f(-k)\)。
\end{solution}

\begin{solution}{fourier-dirichlet-kernel-formula}
\[
 \sum_{k=-N}^{N}e^{ikt}
 =e^{-iNt}\frac{1-e^{i(2N+1)t}}{1-e^{it}}
 =\frac{\sin((N+\tfrac12)t)}{\sin(t/2)}.
\]
在 \(t=0\) 处按连续延拓取值 \(2N+1\)。逐项积分时只有 \(k=0\) 项保留，
故 \((2\pi)^{-1}\int_{-\pi}^{\pi}D_N=1\)。
\end{solution}

\begin{solution}{fourier-translation-coefficients}
周期换元 \(u=x-a\) 给
\[
 \widehat g(k)=\frac1{2\pi}\int f(x-a)e^{-ikx}\,dx
 =e^{-ika}\frac1{2\pi}\int f(u)e^{-iku}\,du.
\]
积分区间平移后仍覆盖一个完整周期。
\end{solution}

\begin{solution}{fourier-step-series}
函数为奇函数，余弦系数全为零。计算得
\[
 b_k=\frac2\pi\int_0^\pi\sin(kx)\,dx
 =\begin{cases}4/(\pi k),&k\text{ 为奇数},\\0,&k\text{ 为偶数}.\end{cases}
\]
所以
\[
 S_Nf(x)=\frac4\pi
 \sum_{\substack{1\le k\le N\\k\ {\rm odd}}}\frac{\sin(kx)}k.
\]
\end{solution}

\begin{solution}{fourier-dini-piecewise-c-one}
若 \(f\) 在 \(x\) 的右侧 \(C^1\)，中值定理给
\(|f(x+t)-f(x+)|\le Ct\)；左侧同理。因此 Dini 商在零附近有界，在
远离零处因 \(f\in L^1\) 可积。定理给
\[
 S_Nf(x)\to\frac{f(x+)+f(x-)}2.
\]
连续点两侧极限均为 \(f(x)\)；跳跃点得到算术平均。
\end{solution}

\begin{solution}{fourier-localization}
令 \(h=f-g\)，则 \(h\) 在 \(x\) 的某邻域为零。由卷积公式，
\[
 S_Nh(x)=\frac1{2\pi}\int_{\delta\le|t|\le\pi}
 h(x-t)\frac{\sin((N+\tfrac12)t)}{\sin(t/2)}\,dt.
\]
被积振幅 \(h(x-t)/\sin(t/2)\) 在该区域属于 \(L^1\)，
Riemann--Lebesgue 引理给极限为零。
\end{solution}

\begin{solution}{fourier-fejer-three-properties}
平方公式直接给 \(K_N\ge0\)。它是 Dirichlet 核平均，故积分平均为一。
若 \(\delta\le|t|\le\pi\)，则
\[
 0\le K_N(t)\le\frac1{(N+1)\sin^2(\delta/2)}\to0.
\]
这也给出该区域积分质量趋零。
\end{solution}

\begin{solution}{fourier-fejer-positive-contraction}
记归一化卷积测度为 \(dt/(2\pi)\)。对 \(p<\infty\)，Jensen 不等式给
\[
 |\sigma_Nf(x)|^p
 \le\frac1{2\pi}\int |f(x-t)|^pK_N(t)\,dt.
\]
再积分并用 Tonelli 与周期换元，得
\(\|\sigma_Nf\|_p^p\le\|f\|_p^p\)。\(p=\infty\) 由核非负且质量一直接得到。
\end{solution}

\begin{solution}{fourier-fejer-lp}
连续周期函数在 \(L^p\) 中稠密。给定 \(f\in L^p\)，取连续 \(g\) 使
\(\|f-g\|_p<\varepsilon\)。收缩性和 Fejér 一致收敛给
\[
 \|\sigma_Nf-f\|_p
 \le2\|f-g\|_p+\|\sigma_Ng-g\|_p
 <2\varepsilon+o(1).
\]
令 \(\varepsilon\downarrow0\) 即得。
\end{solution}

\begin{solution}{fourier-uniqueness-lone}
全部系数为零时，每个 \(\sigma_Nf\) 都为零。上一题给
\(\sigma_Nf\to f\) 于 \(L^1\)，故 \(\|f\|_1=0\)，即 \(f=0\) 几乎处处。
\end{solution}

\begin{solution}{fourier-dirichlet-lone-growth}
由
\[
 |D_N(t)|=\frac{|\sin((N+\tfrac12)t)|}{|\sin(t/2)|}
\]
及 \(c|t|\le|\sin(t/2)|\le C|t|\)（\(0<|t|\le\pi\)），问题化为估计
\(\int_0^\pi|\sin((N+\tfrac12)t)|t^{-1}dt\)。令
\(u=(N+\tfrac12)t\)，再按长度 \(\pi\) 的区间分块。第 \(j\) 块的积分
上下界均为常数乘 \(1/j\)，有限个初始块由常数吸收。因此总和与
\(\sum_{j\le N}1/j\) 同阶，即为 \(\log(N+2)\)。
\end{solution}

\begin{solution}{fourier-gibbs-scaling}
只取奇数 \(k=2m+1\)。当 \(x=y/N\) 时，
\[
 \frac4\pi\sum_{2m+1\le N}\frac{\sin((2m+1)y/N)}{2m+1}
\]
是以步长约 \(2/N\) 对下式积分的 Riemann 和：
\[
 \frac2\pi\int_0^1\frac{\sin(ys)}s\,ds
 =\frac2\pi\int_0^y\frac{\sin u}u\,du.
\]
因此左端趋于右端；在 \(y=\pi\) 处得到正文超调常数。
另一方面，若误差主要位于宽度 \(O(N^{-1})\)、高度有界的邻域，其平方积分
为 \(O(N^{-1})\)，仍趋于零；固定点处也按 Dirichlet--Jordan 定理收敛。
\end{solution}
